% Architecture Worksheet — a copyable companion to "From Token to Fleet"
% Printable one-page-per-question sheet that turns the book's decision loop into
% a fill-in artifact an architect can keep beside the keyboard.

\chapter{The Architecture Worksheet}
\label{chap:worksheet}

\section*{How to use this}

Every chapter of this book turns a vague ask into a defensible number. This
worksheet is the reverse: a single sheet that applies the whole loop to your
own workload. Fill it in row by row; each row is a decision in Chapters
1--26. When a row is empty, that is precisely the gap to investigate next.
The canonical workload box (Chapter 4) shows what a fully-filled row looks
like; the two derived rates most often needed are given under its table.

All quantities are the architect's own --- traceable, verifiable, and stated
in the evidence taxonomy ($[1P]$ first-party, $[2^\circ]$ derived, HYPOTHESIS
testable-but-unmeasured).

\section*{Row 1 --- Requirement} \label{ws:req}

\begin{quote}
\emph{Business objective:} \hrulefill

\emph{Constraints (privacy, region, ops, budget):} \hrulefill
\end{quote}

\section*{Row 2 --- Workload} \label{ws:workload}

\begin{tabular}{@{}p{4.6cm}p{3.0cm}p{5.2cm}@{}}
\toprule
\textbf{Quantity} & \textbf{Measured / target} & \textbf{Evidence} \\
\midrule
Registered users (or seats) & & \\
Concurrency (\% active) & & \\
Average request rate (rps) & & \\
Peak request rate (rps) & & \\
Input tokens / request & & \\
Output tokens / request & & \\
TTFT budget / SLO & & \\
TPOT budget / SLO & & \\
\bottomrule
\end{tabular}

\section*{Row 3 --- Model} \label{ws:model}

\begin{tabular}{@{}p{4.6cm}p{3.0cm}p{5.2cm}@{}}
\toprule
\textbf{Quantity} & \textbf{Value} & \textbf{Evidence} \\
\midrule
Model class (dense / MoE, total params) & & \\
Active params / token & & \\
Weight precision; bytes/param & & \\
Weight residency & & \\
KV precision; KV per token & & \\
KV residency (weights + KV) & & \\
\bottomrule
\end{tabular}

\section*{Row 4 --- Compute \& memory} \label{ws:compute}

\begin{tabular}{@{}p{4.6cm}p{3.0cm}p{5.2cm}@{}}
\toprule
\textbf{Quantity} & \textbf{Derived} & \textbf{Evidence} \\
\midrule
FLOPs / token (prefill, decode) & & \\
Arithmetic intensity (FLOP/byte) & & \\
Ridge point of the hardware & & \\
Compute- or memory-bound? & & \\
MFU (measured) & & \\
\bottomrule
\end{tabular}

\section*{Row 5 --- Architecture implication} \label{ws:arch}

\begin{quote}
\emph{So the architecture does X, and not Y, because:} \hrulefill

\emph{Parallelization strategy:} \hrulefill

\emph{Batching / serving choices:} \hrulefill
\end{quote}

\section*{Row 6 --- Decision record} \label{ws:decision}

\begin{quote}
\emph{Decision:} \hrulefill

\emph{Alternatives rejected (and why):} \hrulefill

\emph{Consequences / monitoring to add:} \hrulefill
\end{quote}

\section*{Following up}

Keep one sheet per significant architecture decision, and re-check it when
the workload's numbers change. The worksheet's sequence --- requirement,
workload, model, compute/memory, architecture, decision --- is the same loop
Chapter 22 makes explicit, and the same reason the Patterns Library in
Chapter 26 exists: a decision recorded with its evidence can be reused, while
a decision recorded without it must be re-derived every time.
